\subsection{Sources of Bias in AI Data and Algorithms}
\label{sec:6.1.1}
Bias in AI systems enters through data collection, data labeling, algorithm design, and implementation, producing unfair, discriminatory, or otherwise undesirable outcomes wherever the system is deployed.

\runin{Data Collection}
An unrepresentative dataset produces bias at the source: a facial recognition system trained mostly on light-skinned faces performs worse on darker skin tones, misidentifying or failing to recognize them at all. In 2018 the ACLU tested Amazon's Rekognition against a public database of arrest photos, running every sitting member of Congress through it. The system returned 28 false matches, a disproportionate share of them non-white members of Congress — a direct result of training data that overrepresented lighter skin tones \autocite{snow2018amazons}.

\runin{Data Labeling}
Human annotators bring their own biases into the labeling process, consciously or not — an image of a woman more likely tagged "nurse," an image of a man tagged "doctor," reinforcing the stereotype in whatever a system trained on those labels later predicts.

Whose judgment is being encoded is a question about working conditions before it is a question about cognition, and this book would be evading something if it discussed labeling bias without saying who does the labeling. The labels that make a system safe are produced by people paid at the bottom of the supply chain to look at the material the system is being taught to refuse. To build the filter that keeps ChatGPT from reproducing the worst of its training data, OpenAI contracted with an outsourcing firm whose workers in Nairobi read and categorized descriptions of child sexual abuse, torture, bestiality, murder, and suicide, for take-home wages between roughly \$1.32 and \$2 an hour. The workers interviewed about it described lasting psychological injury; the counseling they were entitled to was rare and, by their account, useless against production quotas. The contract was terminated eight months early \autocite{perrigo2023openai}.

Two things follow, and the second is the one that gets skipped. The first is the ordinary labor point: this is harmful work, priced as though it were not. The second is epistemic. A judgment made at speed, under a quota, by someone with no stake in the outcome and no channel for saying that the category itself is wrong is a judgment with a known failure mode, and it is the judgment the model learns. An annotation pipeline built so that the annotator cannot dissent produces exactly the ground truth it was built to produce. That is the first feature of section 2.1.4's structural signature, showing up somewhere the political version of the argument never looks: a channel that collects disagreement and is built so that none of it can move anything. Chapter 7 takes this up at length, alongside the same arrangement one stage downstream in how models are tuned on human feedback. The point here is prior to both. A system's values are downstream of a labor arrangement, and the arrangement is a design choice.

The same question runs one stage further back, and stopping at the labels would be the same evasion a second time. Labels are applied to material, and the material was collected. The industry's standard method was to take text and images at web scale without asking, and courts have begun to sort out what that means. A US ruling in Bartz v. Anthropic held that training a model on books the company had bought was transformative and protected as fair use, and that downloading the same books from piracy sites was not; the distinction turned on acquisition alone, and it produced the largest copyright settlement on record \autocite{bartz2025fairuse}. The defendant is named here on purpose. It is the company whose model helped produce this book, disclosed in chapter 0, and a chapter arguing that provenance should be visible cannot describe a provenance case in the passive voice. The European Union now requires the provider of a general-purpose model to publish a summary of what it trained on, against a template that asks for sources by category, scraped web content among them \autocite{europeanunion2024aiact}. Section 8.7.6 records China imposing a disclosure requirement of its own, from a different motive.

What follows here is not the rights question, which is being settled elsewhere. The argument that a web-scale corpus cannot be documented or audited at the size involved, and so encodes a hegemonic slice of the web while presenting as general, was made early and has not been answered by anything that came after \autocite{bender2021parrots}. The point to add here is that a corpus assembled this way has no channel at all. An annotator has a queue and a guideline and could in principle be given a way to object; a writer or photographer whose work sits in a training set was never party to anything, so there is nothing for a channel to attach to. Chapter 7 describes a mechanism that collects disagreement and moves none of it, and section 7.3 states the obligation that follows for a book making this argument. The stage described here sits before both, because no disagreement is collected where none is asked for.

\runin{Algorithm Design}
An algorithm not built to account for bias can encode it structurally rather than merely inherit it from data — penalizing loan applicants for living in a particular neighborhood, say, with a disparate impact on minority and low-income communities regardless of what an applicant's own qualifications show. ProPublica's 2016 investigation into COMPAS, a recidivism-risk tool used in criminal sentencing, found exactly this kind of structural disparity: Black defendants who did not reoffend were nearly twice as likely as white defendants who did not reoffend to be flagged high-risk, while white defendants who did reoffend were scored low-risk more often than Black defendants who did \autocite{angwin2016machine}.

A model does not merely inherit whatever skew its data carries. It can sharpen it. A 2017 study of image datasets used for activity recognition found that cooking scenes in the training set were over 33 percent more likely to depict a woman than a man — an imbalance in the data itself, not a mislabeling of it — and that a model trained on that data amplified the imbalance to 68 percent when tested. The system did not learn the correlation the data showed; it learned a stronger one, because a stronger one predicted better \autocite{zhao2017men}. This is a distinct failure from the one COMPAS illustrates, and it is worse in a specific way: auditing the training data for balance is not sufficient, because the disparity the deployed system acts on can exceed anything visible in the data it was trained on.

\runin{Implementation}
Deployment context can introduce bias that neither the training data nor the algorithm shows on its own. Researchers who fed historical drug-arrest data from Oakland into PredPol's published predictive-policing algorithm in 2016 found it would dispatch patrols to Black neighborhoods at roughly twice the rate of white neighborhoods — a gap survey data on actual drug use across races does not support. The distortion traces back to what the arrest data was measuring: not where drug use was happening, but where police had already been sent to look for it, a disparity the algorithm read as a crime-rate difference and then entrenched further \autocite{lum2016predict}.
